\subsection{The three logically separate parts of the cuspidal criterion}

The following statements separate the construction, its evaluation,
and the exceptional-set input. This separation is essential to the
two different proofs of A.1. In each statement $k$ is finite of
characteristic $p$, $C/k$ is smooth proper geometrically connected
of genus at least two, and $r\ne p$ belongs to the retained prime
set $\Sigma$. The group $\Pi_C^\Sigma$ retains the whole $G_k$.
For $A\in C(k)$, $G_A$ is the arithmetic maximal cuspidally
pro-$r$ quotient of $C-A$ relative to $\Pi_C^\Sigma$.

\begin{claim}{A.42}{Proved}{Central extensions and the difference cocycle}
Suppose $P,Q\in C(k)$ and
$\beta:G_P\xrightarrow{\sim}G_Q$ is a continuous isomorphism over
$\operatorname{id}_{\Pi_C^\Sigma}$ preserving the cusp decomposition
conjugacy class. Put $M=\mathbf Z_r(1)$, with its arithmetic
cyclotomic action and its actual local inertia convention:
$z^{1/r^a}\mapsto\zeta z^{1/r^a}$ represents $\zeta\in\mu_{r^a}$.
The geometrically cuspidally central quotients
fit into exact sequences
\[
 1\to M\to E_A\to\Pi_C^\Sigma\to1\qquad(A=P,Q).
\]
The isomorphism induced by $\beta$ is the identity on $M$.
On the arithmetic geometrically pro-$\Sigma$ group $\Gamma$ of
$C-\{P,Q\}$ there is a continuous $M$-valued crossed homomorphism
$\chi$ with residues $(1,-1)$ at $(P,Q)$. If
$\operatorname{div}(f)=n(P-Q)$ for $f\in k(C)^\times$ and $n>0$,
then there is $a\in k^\times$ such that
\[
                  \kappa_r(f)-n[\chi]=\kappa_r(a)
                  \quad\text{in }H^1(\Gamma,M).           \tag{A.42.1}
\]
No assertion about exceptional points or transported point labels
is used in this construction.
\end{claim}
\begin{proof}
Quotient the cuspidal kernel by its commutators with the entire
geometric subgroup. The remaining kernel is a quotient of $M$:
the cusp kernel is normally generated by inertia at the single
puncture, and all its conjugates now agree. Its arithmetic action
is cyclotomic, as is the action on tame inertia.

We check that no further quotient of $M$ has occurred by an
algebraic argument. For each $a\ge1$, the degree isomorphism is
\[
 H^2(C_{\bar k},\mu_{r^a})\simeq\mathbf Z/r^a,
                    \qquad c_1(\mathcal O_C(A))\longmapsto1.
\]
This is the smooth proper curve case of
\citet[Proposition 14.2 and Theorem 24.1(a)]{MilneEtale}.
The comparison with $H^2(\Delta_C^\Sigma,\mu_{r^a})$ follows
without lifting the curve: in the tower of finite geometric
$\Sigma$-covers, a degree-one cohomology class is killed by its
finite cyclic torsor, and a degree-two class is killed by an
\emph{\'etale} cover of degree divisible by $r^a$. Such covers
exist because a positive-genus Jacobian has nonzero prime-to-$p$
Tate module and the Abel map identifies its abelian covers.
The same argument applies after each cover; cohomology above
degree two is zero. The Hochschild--Serre tower argument in
\citet[Theorem 14.9 and Remark 14.10]{MilneEtale} therefore gives
the comparison. The arithmetic comparison follows by applying
Hochschild--Serre for the unchanged $G_k$ to both sides.

Consider the gerbe of $r^a$-th roots of $\mathcal O_C(A)$.
Its canonical trivialization on $C-A$ gives a homomorphism from
the punctured group to the corresponding extension of
$\Pi_C^\Sigma$ by $\mu_{r^a}$. Fix the band identification once,
at every compatible level, by requiring that the local Kummer
coordinate $w^{r^a}=z$ identify the image of actual inertia with
$w\mapsto\zeta w$, hence with $\zeta$. Its extension class is
$\eta_{A,a}=\varepsilon c_1(\mathcal O_C(A))$, where the single
sign $\varepsilon\in\{1,-1\}$ depends only on the convention
relating gerbes and group-extension factor sets. It is independent
of $A$ and of $a$: the local root chart and its band identification
are the same and commute with the power transition maps. This
fixes the local residue at $+1$ without requiring a convention for
the sign of a spectral-sequence transgression. The homomorphism
is surjective:
the punctured group surjects onto the proper group, and its
inertia surjects onto the kernel. It factors through the
geometrically central cuspidal quotient, since its geometric
kernel is central. Thus that quotient retains inertia modulo
every $r^a$. A closed subgroup of $\mathbf Z_r$ is either zero
or $r^b\mathbf Z_r$; it follows that its inertia kernel is zero.
This proves the exact sequence with full kernel $M$ and also
identifies the geometric extension class with the primitive
orientation class $\eta_{A,a}$ normalized by that actual inertia.
For any two rational points $P,Q$, the degree isomorphism gives
$\eta_{P,a}=\eta_{Q,a}$, and both are generators. All these arguments work in characteristic
$p$, including $p=2$ and $r=2$ when $r\ne p$.

The induced map $\bar\beta:E_P\to E_Q$ acts on the procyclic
kernel by a unit $u\in\mathbf Z_r^\times$. Naturality of
extension classes over the identity of the geometric base gives
$u\eta_{P,a}=\eta_{Q,a}=\eta_{P,a}$. Since this common class
is primitive, $u\equiv1$
modulo every $r^a$, and $u=1$.

On the continuous fibre product $E_P\times_{\Pi_C^\Sigma}E_Q$
put $\chi(e_P,e_Q)=\bar\beta(e_P)e_Q^{-1}\in M$, using additive
notation on $M$. Direct multiplication gives
$\chi(ee')=\chi(e)+e\chi(e')$. Pull back along the two natural
punctured-group maps from $\Gamma$. At $P$ the second map kills
inertia and the first has residue $1$; at $Q$ the first kills
inertia and the inverse second map has residue $-1$.
The Kummer cocycle of $f$ has residues $(n,-n)$, so their
difference vanishes on both inertia groups. The coefficient
module is trivial under the geometric kernel; the cocycle
therefore vanishes on their closed normal closure. That closure
is the kernel of the proper quotient. The difference descends
to $H^1(\Pi_C^\Sigma,M)$.

Finally inflation is an isomorphism
\[
                  H^1(G_k,M)\xrightarrow{\sim}
                  H^1(\Pi_C^\Sigma,M).
\]
Indeed the possible extra term is
$H^1(\Delta_C^\Sigma,M)^{G_k}$. This is a submodule of
$\operatorname{Hom}(T_rJ_C,\mathbf Z_r(1))$, whose Frobenius
eigenvalues have complex absolute value $\sqrt{\#k}$; none is
$1$. The module is torsion-free, so its invariant submodule
vanishes. Finite-field Kummer theory identifies the left side
with the $r$-primary completion of the finite group $k^\times$.
It is represented by some $a\in k^\times$, proving (A.42.1).
\end{proof}

\begin{claim}{A.43}{Proved}{Evaluation at a point with unchanged label}
Keep the hypotheses and notation of A.42. Let $R$ be a closed
point different from $P,Q$, put $L=k(R)$, and suppose that
$\beta$ carries a point decomposition subgroup at $R$ in $G_P$
to a geometric conjugate of the point decomposition subgroup
at that same $R$ in $G_Q$. Then
\[
             [\chi]|_{s_R(G_L)}=0,
             \qquad \kappa_{r,L}(f(R)/a)=0.               \tag{A.43.1}
\]
There is no finiteness hypothesis on the set of points where
this same-label assumption holds.
\end{claim}
\begin{proof}
Choose the two point subgroups to have the same image $D_R$
in $\Pi_{C_L}^\Sigma$. A conjugator in the target has image
normalizing $D_R$. A.17, equivalently
\citet[Lemma 1.7]{ST2009}, says its normalizer is exactly
$D_R$: the point subgroup projects isomorphically onto the
entire $G_L$. Multiply the conjugator on the appropriate side
by the unique element of the target point subgroup with
that same proper image. Its conjugation of the point subgroup
does not change, and its proper image becomes $1$.
In $E_Q$ it therefore belongs to $M$. Conjugating a splitting
by an element of $M$ changes its $M$-cocycle by a coboundary.
The difference cocycle consequently restricts to zero in
$H^1(G_L,M)$. Restrict (A.42.1) along the actual point section;
functorial Kummer theory evaluates $f$ at $R$, yielding the
second equality. The field extension is the residue field of
this individual point; no common field for all $R$ is required.
\end{proof}

\begin{claim}{A.44}{Proved}{The finite-exception endgame}
Keep the hypotheses of A.42. If all but finitely many closed
points of $C$ satisfy the same-label assumption of A.43, then
$P=Q$.
\end{claim}
\begin{proof}
Assume $P\ne Q$. Since $J_C(k)$ is finite, the class $[P-Q]$
has finite positive order $n$. A geometric function with divisor
$n(P-Q)$ descends to $k(C)$: its Galois ratios are constants,
and Hilbert 90 removes their scalar cocycle. Thus choose $f$
as in A.42. The finite excluded set has only finitely many
geometric points and hence finitely many values under $f/a$.
The finite nonconstant map $f/a:C_{\bar k}\to\mathbf P^1$
is surjective. There are infinitely many $r$-power roots of
unity in $\bar k$, because $r\ne p$. Choose a nonidentity
one, $\zeta$, outside the finite excluded values, and a
geometric point above it. Its residue field $L$ is finite,
and its underlying closed point satisfies A.43. But $f(R)/a$
is the nonzero $r$-primary element $\zeta$ of $L^\times$,
whose image in its $r$-primary completion is nonzero.
This contradicts (A.43.1). Hence $P=Q$.
\end{proof}

\paragraph{Check record and exact dependency separation.}
Checked only in this writing as separate statements: GPT-6 Astra
(OpenAI, ChatGPT Pro mode), 18 November 2026, direct proof
checking and algebraic rederivation. They replace the corresponding
three portions of A.19; those portions were also checked in this
writing after the earlier readings. A.42 depends on curve
cohomology, Abel covers and Kummer theory, and A.43 additionally
uses A.17. Neither uses A.16. A.44 is where finiteness of the
exceptional labels is used. In the alternative proof, A.20
uses A.42--A.43 and the Weil bounds, while A.21 supplies
proper-field support by the degree-two-pencil argument and
the complementary torsion representation. It uses neither
A.16 nor A.44. Thus the dependency edge A.20$\to$A.19
is replaced by A.20$\to$A.42,A.43; the edge through A.16
is removed rather than hidden in a citation to the earlier
unsplit criterion.
